\emph{Preliminaries.} By Theorem \claimref{tower.ladder-pairs-satisfy-car-vectors} each plane gives a mode $a^{*} = v$, $a = w$ with $\{a_{x}, a_{y}^{*}\} = \delta_{xy}$, $\{a_{x}, a_{y}\} = 0$, and $v^{*} = w$. One plane's algebra is spanned by $1, v, w, vw$ (since $e = v + w$, $f = v - w$, $ef = wv - vw$), so $\mathcal{A}$ is spanned by the ordered monomials $\prod_{x}m_{x}$ with $m_{x} \in \{1, a_{x}^{*}, a_{x}, a_{x}^{*}a_{x}\}$, and these are a basis (there are $4^{n}$ of them for $n$ planes, the dimension of the algebra). Call a monomial's \emph{doublet degree} the number of its factors $m_{x} \in \{a_{x}^{*}, a_{x}\}$ with $x$ a doublet channel, and likewise its singlet degree; their parities are $p_{D}$, $p_{S}$. Writing a monomial with its doublet factors first (a sign) shows it is a product $xy$ with $x \in \mathcal{A}_{D}$ of parity $p_{D}$ and $y \in \mathcal{A}_{S}$ of parity $p_{S}$, and conversely every such product is a combination of monomials with these parities. So $\mathcal{A} = \bigoplus\mathcal{A}^{(p_{D}, p_{S})}$.

\emph{(i).} The isospin generators and $Y$ are combinations of $a_{x}^{*}a_{y}$, even in both kinds of modes. From the relations, $[a_{x}^{*}a_{y}, a_{z}^{*}] = \delta_{yz}a_{x}^{*}$ and $[a_{x}^{*}a_{y}, a_{z}] = -\delta_{xz}a_{y}$, so $\operatorname{ad}(a_{x}^{*}a_{y})$ is a derivation that maps each doublet mode operator to a doublet mode operator (or zero) when $x, y$ are doublet channels, and each singlet one to a singlet one when they are singlet channels; hence $\operatorname{ad}T_{j}$ and $\operatorname{ad}Y$ preserve the doublet and singlet degrees of every monomial up to terms of the same parities, and each $\mathcal{A}^{(p_{D}, p_{S})}$ is invariant. In particular $[T_{3}, a_{\delta\uparrow}^{*}] = \tfrac12a_{\delta\uparrow}^{*}$, $[T_{3}, a_{\delta\uparrow}] = -\tfrac12a_{\delta\uparrow}$, $[T_{3}, a_{\delta\downarrow}^{*}] = -\tfrac12a_{\delta\downarrow}^{*}$, $[T_{3}, a_{\delta\downarrow}] = \tfrac12a_{\delta\downarrow}$, and $\operatorname{ad}T_{3}$ kills $a_{x}^{*}a_{x}$ and the singlet modes; as a derivation it is diagonal on the monomials, with eigenvalue the sum of the factors' weights, an element of $\tfrac12(\text{doublet degree}) + \mathbb{Z} = \tfrac12p_{D} + \mathbb{Z}$. An element with a component of half-integer isospin has a component with a half-integer eigenvalue of $\operatorname{ad}T_{3}$, so its $\mathcal{A}^{(0, \cdot)}$ part cannot carry it: every $\operatorname{ad}T_{3}$-eigenvector with half-integer eigenvalue lies in $\mathcal{A}^{(1, 0)} \oplus \mathcal{A}^{(1, 1)}$, and if it is even (total parity $p_{D} + p_{S}$ even) it lies in $\mathcal{A}^{(1, 1)}$, which is zero when $S$ is empty. For one doublet pair of modes, $\mathcal{A}_{D}$ is the algebra of the two rung-type planes of Theorem \claimref{synthesis.rung-field-is-odd-in-the-tower-and-its-modes-carry-no-even-doublet}(2): its odd part (dimension $8$) has Casimir $\tfrac34$ and its even part (dimension $8$) Casimir eigenvalues $0$ (multiplicity $5$) and $2$ (multiplicity $3$). Since $\operatorname{ad}T_{j}$ acts on $xy$ as $(\operatorname{ad}T_{j}x)y$ ($T_{j}$ commutes with $\mathcal{A}_{S}$, being even and built from doublet modes, which anticommute with singlet modes), $\mathcal{A}^{(p_{D}, p_{S})} \cong \mathcal{A}_{D}^{p_{D}}\otimes\mathcal{A}_{S}^{p_{S}}$ as an isospin representation with $\mathcal{A}_{S}$ inert, which gives the stated Casimir values and multiplicities. For the bilinears, take $\delta$ fixed and compute with the two rules above: $[T_{+}, a_{\delta\uparrow}^{*}a_{s}] = 0$, $[T_{-}, a_{\delta\uparrow}^{*}a_{s}] = a_{\delta\downarrow}^{*}a_{s}$, $[T_{+}, a_{\delta\downarrow}^{*}a_{s}] = a_{\delta\uparrow}^{*}a_{s}$ (with $T_{+} = \sum a_{\delta\uparrow}^{*}a_{\delta\downarrow}$), and $T_{3}$ weights $\pm\tfrac12$: a doublet. The other three pairs follow in the same way (the minus signs in the lower components come from $[T_{-}, a_{\delta\downarrow}] = -a_{\delta\uparrow}$, with $T_{-} = \sum a_{\delta\downarrow}^{*}a_{\delta\uparrow}$). For hypercharge, $[Y, a_{x}^{*}] = Y_{x}a_{x}^{*}$ and $[Y, a_{x}] = -Y_{x}a_{x}$ give the four values listed.

\emph{(ii).} A channel matrix $M$ with $[W, M] = hM$ has nonzero entries only where $w_{\alpha} - w_{\beta} = h$; a nonzero entry of $M^{n}$ needs a chain with weight differences summing to $nh$, impossible once $n|h|$ exceeds the spread of the weights, so $M$ is nilpotent, and a normal nilpotent matrix is zero. This is the argument of Theorem \claimref{synthesis.helicity-one-pair-modes-on-a-null-generator-are-quasi-spins}(iii), which uses only that the weights are real numbers; the weights there are the helicity weights, here the hypercharge or isospin weights. By its clause (ii) a nonzero non-normal $M$ gives a $\rho_{q}(M)$ that is not a canonical boson on the edge-clean vectors ($N > q$), and by its clause (iv), when $M^{2} = 0$, $B = \rho_{q}(M)^{*}$ on the chiral sea has $|B^{p}\Omega|^{2} = (p!)^{2}e_{p}$ over the $q\operatorname{rank}M$ numbers $\sigma_{s}^{2}$, zero for $p > q\operatorname{rank}M$. A component of a doublet of pair densities has $T_{3} = \pm\tfrac12$: its matrix has $[W_{3}, M] = \pm\tfrac12M$ for the isospin weights $W_{3}$ ($\pm\tfrac12$ on the members of each doublet, $0$ on singlets); and if the doublet has hypercharge $y \ne 0$, also $[W_{Y}, M] = yM$. For $M = E_{(\delta\uparrow), s}$, $M^{2} = 0$ and $\operatorname{rank}M = 1$, the singular value is $1$ and $e_{p}$ of $q$ ones is $\binom{q}{p}$. On edge-clean vectors the window identity of the cited clause (i) gives the value $q\operatorname{tr}(M^{*}M) + \rho_{0}([M, M^{*}]) = q + \rho_{0}(E_{\delta\uparrow,\delta\uparrow} - E_{ss})$, which is $q + 1$ with one extra particle in channel $\delta\uparrow$ at a bulk momentum and $q - 1$ with one in channel $s$; and for $M' = E_{(\delta\downarrow), s}$ it gives $[\rho_{q}(M), \rho_{q}(M')^{*}] = q\operatorname{tr}(M'^{*}M) + \rho_{0}([M, M'^{*}]) = \rho_{0}(E_{\delta\uparrow,\delta\downarrow})$, since $M'^{*}M = E_{s,\delta\downarrow}E_{\delta\uparrow,s} = 0$ and $MM'^{*} = E_{\delta\uparrow,\delta\downarrow}$; on an edge-clean vector the edge momenta are filled or empty in every channel, where $a_{k,\delta\uparrow}^{*}a_{k,\delta\downarrow}$ acts as zero, so only the bulk remains.

\emph{Reading (argued).} The rung field's oddness (Theorem \claimref{synthesis.rung-field-is-odd-in-the-tower-and-its-modes-carry-no-even-doublet}) does not exclude a composite doublet of two kinds of fermions: (i) says exactly what it requires, a singlet constituent, and how the hypercharges must match. Particle-particle doublets reduce to the hopping case by the particle-hole reflection of the singlet planes ($f \mapsto -f$ on those planes exchanges $v_{s}$ and $w_{s}$, reverses the singlet sea and the sign of $Y_{s}$), so they too are nilpotent with charge $\mp(Y_{\delta} + Y_{s})$ and isospin $\pm\tfrac12$. So no component of a composite doublet is among the exact (normal, neutral) pair bosons of a generator; each is a quasi-spin of finite capacity, canonical only in a contraction over many sites, as the helicity-one modes are for the photon. This bears on the boson statistics of a composite Higgs, not on its gap.
